%% kapa: exam=1 round=37 year=2026 period=1 form=A subject=economics subject_ko=경제학원론 seq=17 number=57
%% source: Q-Net 감정평가사 기출문제 (cst00309) artlSeq=5258924 file=raw/1008/5258924/2026년 제37회 감정평가사 1차 1교시.pdf
%% conversion: kordoc 4.24.51 -> tools/gen1.py (auto); review=proofread-sonnet; flags=none
\begin{question}{57}
외부성에 관한 설명으로 옳지 않은 것을 모두 고른 것은?
\begin{framedtext}
ㄱ. 양(+)의 외부성이 존재하면 시장 생산량은 사회적 최적 생산량보다 많다. \par
ㄴ. 외부성은 보조금 혹은 조세 등을 통해 내부화시킬 수 있다. \par
ㄷ. 코우즈(R. Coase)에 의하면 외부성이 존재하더라도 재산권이 확립되어 있으면 당사자 간 협상으로 외부성을 치유할 수 있다. \par
ㄹ. 외부성의 치유방법 중 정부의 직접규제가 사회 전체적인 비용이 가장 작다. \par
\end{framedtext}
\begin{choices}
  \item ㄱ, ㄴ
  \item ㄱ, ㄹ
  \item ㄴ, ㄷ
  \item ㄷ, ㄹ
  \item ㄱ, ㄴ, ㄷ, ㄹ
\end{choices}
\end{question}
\answer{2}
